% Exercise 3.13

\begin{itemize}
    \item[(a)] We have the following \verb|R| output:
        \scriptsize\begin{verbatim}
> x <- rnorm(100, 0, 1)
> x
  [1]  1.134965089  1.111931845 -0.870777634  0.210731585  0.069395647
  [6] -1.662648853  0.810839980 -1.912345796 -1.246753429  0.998154445
 [11] -0.540872745 -0.216375791 -1.621937293 -1.450963965  0.350909731
 [16] -0.174546929 -0.591428470 -1.334027261 -1.097298501  2.036103609
 [21] -0.326489593  0.774005212  0.785006401  0.763246080  0.294808760
 [26] -1.252355924 -1.009503753  0.751391195 -1.308353513  0.527540097
 [31] -0.533539574 -0.398376014 -0.789569450 -0.230141136  0.877184842
 [36]  0.453733178 -0.232464148  0.870005525  1.656003734 -0.006368929
 [41]  0.470489453  0.278218649 -0.977902941 -0.926586142  1.919770463
 [46]  0.881277788  0.742081772  0.147573404  0.485388565  0.151856040
 [51]  0.041998754  0.223422312 -1.010465086  2.401222102  0.801961790
 [56] -0.251207959  1.212889371 -0.627258086  1.711158507 -0.394373553
 [61] -2.321490856  1.364119195  1.132229133 -0.774316319 -1.410374966
 [66] -1.834527581 -0.269013538 -1.833928577 -0.814468019  0.163572122
 [71]  0.855519222 -0.819963127 -0.123602760  0.254948236  1.718926338
 [76] -0.958543528 -1.604310262 -1.845609422  0.555737185 -0.060119191
 [81]  0.772086304 -0.140839387  0.393093926  0.224218574  0.023541985
 [86] -0.622962660  1.262009381 -0.405774043  0.666763771  0.164639155
 [91]  1.781524475  0.711213964 -0.337691156 -0.009148952 -0.125309208
 [96] -2.090846097  1.697393895  1.063881154 -0.766616636  0.382007559
        \end{verbatim}\normalsize
    \item[(b)] We have the following output, where the standard deviation
        $\sigma = \sqrt{V(\varepsilon)} = \sqrt{0.25} = 0.0625$:
        \scriptsize\begin{verbatim}
> eps <- rnorm(100, 0, 0.0625)
> eps
  [1] -0.0973147340  0.1201977283 -0.1160518518 -0.1316324022  0.0436030329
  [6]  0.0567152755 -0.0122492624 -0.0129262805  0.0453151982  0.0874199347
 [11] -0.0994096968  0.0815310673  0.0122507334 -0.0212089653  0.0790715088
 [16]  0.0587358553  0.0798720304 -0.0181975090  0.0485107781  0.0184828521
 [21] -0.0323249722  0.1098270640  0.0101079262 -0.0034177735  0.0330607250
 [26]  0.0243584076 -0.0441107134 -0.0033236547  0.1629027805 -0.0045370690
 [31] -0.0185169932  0.0458956460  0.0250925467 -0.0336438173 -0.0676701757
 [36]  0.0529875880  0.0074278134  0.0301404514 -0.0182078003 -0.1032568920
 [41]  0.0043184124  0.1494110437 -0.0913607304 -0.0615503933 -0.0835033590
 [46] -0.0228207257  0.0863408905 -0.0095823285 -0.0159732271 -0.0805391940
 [51]  0.0040791511  0.0647079013  0.1412634870  0.0821685172 -0.0543764591
 [56] -0.0314456272  0.0379659722 -0.0006275885  0.0172081437 -0.0820186740
 [61] -0.0266581167 -0.0579002655 -0.0580348703 -0.0449312876 -0.0259244106
 [66] -0.0003542433  0.0395905575 -0.0318237984 -0.0535630387  0.1010104736
 [71]  0.0621166504  0.0435528662  0.1062481671 -0.0611088716  0.1280510989
 [76]  0.0766530828  0.0191889382  0.0390208115  0.0038338402 -0.0069266120
 [81] -0.0985145975  0.0464118654  0.1338952124  0.1306108098  0.0106121454
 [86] -0.0067423618  0.0113746559  0.0716205111  0.0923311764  0.0254965095
 [91]  0.0696424210 -0.0016922402  0.0311076377  0.0741025009  0.2274733517
 [96] -0.0033751622 -0.0417592073  0.0278629973 -0.0253640779  0.0397051965
        \end{verbatim}\normalsize
    \item[(c)] We have the following \verb|R| commands to produce the vector
        \verb|y|:
        \scriptsize\begin{verbatim}
> y <- -1.0 + 0.5*x + eps
> y
  [1] -0.52983219 -0.32383635 -1.55144067 -1.02626661 -0.92169914 -1.77460915
  [7] -0.60682927 -1.96909918 -1.57806152 -0.41350284 -1.36984607 -1.02665683
 [13] -1.79871791 -1.74669095 -0.74547363 -1.02853761 -1.21584220 -1.68521114
 [19] -1.50013847  0.03653466 -1.19556977 -0.50317033 -0.59738887 -0.62179473
 [25] -0.81953490 -1.60181955 -1.54886259 -0.62762806 -1.49127398 -0.74076702
 [31] -1.28528678 -1.15329236 -1.36969218 -1.14871439 -0.62907775 -0.72014582
 [37] -1.10880426 -0.53485679 -0.19020593 -1.10644136 -0.76043686 -0.71147963
 [43] -1.58031220 -1.52484346 -0.12361813 -0.58218183 -0.54261822 -0.93579563
 [49] -0.77327894 -1.00461117 -0.97492147 -0.82358094 -1.36396906  0.28277957
 [55] -0.65339556 -1.15704961 -0.35558934 -1.31425663 -0.12721260 -1.27920545
 [61] -2.18740354 -0.37584067 -0.49192030 -1.43208945 -1.73111189 -1.91761803
 [67] -1.09491621 -1.94878809 -1.46079705 -0.81720347 -0.51012374 -1.36642870
 [73] -0.95555321 -0.93363475 -0.01248573 -1.40261868 -1.78296619 -1.88378390
 [79] -0.71829757 -1.03698621 -0.71247145 -1.02400783 -0.66955782 -0.75727990
 [85] -0.97761686 -1.31822369 -0.35762065 -1.13126651 -0.57428694 -0.89218391
 [91] -0.03959534 -0.64608526 -1.13773794 -0.93047198 -0.83518125 -2.04879821
 [97] -0.19306226 -0.44019643 -1.40867240 -0.76929102
        \end{verbatim}\normalsize
        The length of \verb|y| is clearly 100 and the values of the 
        coefficients in this linear model are given by $\beta_0 = -1$
        and $\beta_1 = 0.5$.
    \item[(d)]
        In Figure~\ref{fig3_13scat}, we have a scatter plot of the linear model
        $Y = -1 + 0.5X + \varepsilon$.
        \begin{figure}[!ht]
            \includegraphics[scale=0.6, center]{../plots/ex3_13_d.pdf}
            \caption{Scatter plot for $Y = -1 + 0.5X + \varepsilon$ \label{fig3_13scat}}
        \end{figure}
        The scatter plot is quite clearly linear, which makes sense given the
        simple linear model used.
    \item[(e)] We have the following \verb|R| output:
        \scriptsize\begin{verbatim}
> lm.fit <- lm(y ~ x)
> summary(lm.fit)

Call:
lm(formula = y ~ x)

Residuals:
      Min        1Q    Median        3Q       Max 
-0.147641 -0.039818 -0.007492  0.043777  0.212754 

Coefficients:
             Estimate Std. Error t value Pr(>|t|)    
(Intercept) -0.984800   0.006763 -145.62   <2e-16 ***
x            0.503837   0.006600   76.34   <2e-16 ***
---
Signif. codes:  0 ‘***’ 0.001 ‘**’ 0.01 ‘*’ 0.05 ‘.’ 0.1 ‘ ’ 1

Residual standard error: 0.06762 on 98 degrees of freedom
Multiple R-squared:  0.9835,    Adjusted R-squared:  0.9833 
F-statistic:  5827 on 1 and 98 DF,  p-value: < 2.2e-16
        \end{verbatim}\normalsize
        The coefficients in the linear regression are quite close to the actual values, 
        i.e., $\hat{\beta_0} = -0.984800 \approx \beta_0 = -1$ and
        $\hat{\beta_1} = 0.503837 \approx \beta_1 = 0.5$. This fits well with the 
        relatively low $p$-values of $p < 2.2 \times 10^{-16}$.
    \item[(f)]
        In Figure~\ref{fig3_13lines}, we have a scatter plot of the linear model
        with least squares line and population regression line.
        \begin{figure}[!ht]
            \includegraphics[scale=0.6, center]{../plots/ex3_13_f.pdf}
            \caption{Scatter plot with regression lines for $Y = -1 + 0.5X + \varepsilon$ 
                \label{fig3_13lines}}
        \end{figure}
    \item[(g)] We have the following \verb|R| output:
        \scriptsize\begin{verbatim}
> lm.fit <- lm(y ~ x + I(x^2))
> summary(lm.fit)

Call:
lm(formula = y ~ x + I(x^2))

Residuals:
      Min        1Q    Median        3Q       Max 
-0.150041 -0.036205 -0.005171  0.042029  0.210216 

Coefficients:
             Estimate Std. Error  t value Pr(>|t|)    
(Intercept) -0.982249   0.008785 -111.808   <2e-16 ***
x            0.503634   0.006642   75.827   <2e-16 ***
I(x^2)      -0.002432   0.005314   -0.458    0.648    
---
Signif. codes:  0 ‘***’ 0.001 ‘**’ 0.01 ‘*’ 0.05 ‘.’ 0.1 ‘ ’ 1

Residual standard error: 0.0679 on 97 degrees of freedom
Multiple R-squared:  0.9835,    Adjusted R-squared:  0.9832 
F-statistic:  2890 on 2 and 97 DF,  p-value: < 2.2e-16
        \end{verbatim}\normalsize
        Given the high $p$-value of 0.648 associated with $\hat{\beta_2}$ in the model
        $\hat{Y} = \hat{\beta_0} + \hat{\beta_1}X + \hat{\beta_2}X^2$, we see that the
        quadratic term does not improve the model fit, which is to be expected given
        the linear relationship of the underlying data.
\end{itemize}
